\Askhsh{38805}{4}
\begin{ylh}{2}
    \item 2.2. Διάταξη Πραγματικών Αριθμών
    \item 2.3. Απόλυτη Τιμή Πραγματικού Αριθμού
    \item 6.3. Η Συνάρτηση $f(x)=ax+\beta$
\end{ylh}
Το πιο ευρέως χρησιμοποιούμενο υλικό για τη μεταφορά ηλεκτρικής ενέργειας στις μέρες μας είναι ο χαλκός, ο οποίος όχι μόνο είναι εξαιρετικός αγωγός του ηλεκτρισμού, αλλά και ένα φθηνό υλικό με υψηλή ολκιμότητα (μορφοποιείται εύκολα σε σύρμα). Η αντίσταση $(R)$ του χαλκού αλλάζει όταν μεταβάλλεται η θερμοκρασία $(T)$, και συγκεκριμένα αυξάνεται με την αύξηση της θερμοκρασίας. Μέσα σε ένα ευρύ φάσμα θερμοκρασιών (περίπου $- 200{^\circ}C$ έως $500{^\circ}C$), αυτή η σχέση είναι γραμμική, δηλαδή της μορφής $R = a T + \beta$. Για ένα συγκεκριμένο κομμάτι χάλκινου σύρματος, η αντίσταση είναι $R = 12,56\ \Omega\ (Ohm)$ σε θερμοκρασία $T = 0{^\circ}C$. Όταν η θερμοκρασία αυξάνεται στους $27{^\circ}C$, η αντίσταση αυξάνεται στα $13,8\ \Omega$.

\begin{alist}
\item Να βρείτε την εξίσωση που χρησιμοποιείται για τον υπολογισμό της αντίστασης $R$ στο σύρμα σε σχέση με τη θερμοκρασία του $T$.
\monades{7}
\item Να προσδιορίσετε την κλίση της ευθείας $R = a T + \beta$ που βρήκατε στο α) ερώτημα και να επιλέξετε ποια από τις παρακάτω προτάσεις περιγράφει τι αντιπροσωπεύει αυτή η κλίση. Να αιτιολογήσετε την απάντησή σας.

\begin{rlist}
\item Τη μεταβολή της αντίστασης, όταν η θερμοκρασία αυξάνεται κατά $1{^\circ}C$.

\item Τη αντίσταση, όταν η θερμοκρασία είναι $0{^\circ}C$.

\item Τη θερμοκρασία, στην οποία η αντίσταση θα ήταν $0\ \Omega$.

\item Τη μεταβολή της θερμοκρασίας, όταν η αντίσταση αυξάνεται κατά $1\ \Omega$.
\monades{6}
\end{rlist}
\item Γενικά, όσο μικρότερη είναι η αντίσταση στο σύρμα, τόσο μειώνονται οι απώλειες ισχύος, όταν περνάει ρεύμα. Γι' αυτό θέλουμε να έχουμε όσο το δυνατόν μικρότερη αντίσταση. Με δεδομένο ότι σε ένα ευρύ φάσμα θερμοκρασιών η σχέση μεταξύ της αντίστασης και της θερμοκρασίας περιγράφεται από την εξίσωση $R = 0,046T + 12,56$, να βρείτε:

\begin{rlist}
\item Σε ποια θερμοκρασία η αντίσταση θα μειωνόταν στα $10,26\ \Omega$.
\monades{4}
\item Μεταξύ ποιων τιμών θα κυμαίνεται η αντίσταση του σύρματος, αν $d(T, - 50\degree) \leq 50\degree$.
\monades{8}
Δίνεται ότι $1,24 \div 27 \cong 0,046$.

{\footnotesize Το παραπάνω θέμα αναπτύχθηκε στο πλαίσιο του έργου: «Ανάπτυξη Δοκιμασιών Αξιολόγησης Δεξιοτήτων Εγγραμματισμού στα μαθήματα της Νεοελληνικής Γλώσσας και Λογοτεχνίας, της Άλγεβρας, της Φυσικής και της Χημείας Α' Λυκείου Γενικού Λυκείου» Ανάδοχος: «Ειδικός Λογαριασμός Κονδυλίων Έρευνας (Ε.Λ.Κ.Ε) Πανεπιστημίου Ιωαννίνων» ΑΔΑΜ: 25SYMV016348911 2025-02-20.}
\end{rlist}
\end{alist}

